\section{Exact Synthesis of Single Qubit Clifford+T Circuits}

In this section, we present results from Kliuchnikov et al \cite{exact-clifford+t} on an algorithm for the exact synthesis of single-qubit unitaries over the Clifford + T gate set, when it is possible.

The key insight is the following theorem, which relates the set of unitaries that can be implemented by Clifford+T circuits with the ring $\Z\left[\frac{1}{\sqrt{2}}, i\right]$.
\begin{theorem}\label{unitary-iff-clifford+t}
	Denote $\mathcal{CT}(2)$ as the set of all $2 \times 2$ unitary matrices that can be exactly implemented by a single qubit Clifford+T circuit. 
	Then $U \in \mathcal{CT}(2)$ if and only if $U \in U(\Z\left[\frac{1}{\sqrt{2}}, i\right])$, which is to say it is an unitary matrix with entries in $\Z\left[\frac{1}{\sqrt{2}}, i\right]$.

	In such a case, moreover, it can be written in the following form
	\begin{equation}\label{clifford+t-form}
	\begin{pmatrix}
		a & b^\dagger \omega^k \\ 
		b & a^\dagger \omega^k \\	
	\end{pmatrix}
	\end{equation}
	where $a, b \in \Z\left[\frac{1}{\sqrt{2}}, i\right]$, $\omega = e^{\pi i/4}$ and $k$ is an integer, and $\dagger$ denotes the complex conjugate.
\end{theorem}


\begin{remark}\label{rk:first-column-matters}
Notice that 
$$
\begin{pmatrix}
		a & b^\dagger \omega^k \\ 
		b & a^\dagger \omega^k \\	
	\end{pmatrix} T^n 
	= \begin{pmatrix}
		a & b^\dagger \omega^{k+n} \\ 
		b & a^\dagger \omega^{k+n} \\	
\end{pmatrix}
$$
This means we only need to regard the first column of any unitary matrix in finding a Clifford+T circuit that implements it.
As for any $U, V \in \mathcal{CT}(2)$ with the same first column, we have $U = VT^n$ for some integer $n$.
\end{remark}


\begin{definition}[Least Denominator Exponent]\label{def:lde}
	The least denominator exponent of a vector or matrix $u$ with entries in the ring $\Z[\sq, i]$, denoted as $\lde(u)$, is the smallest non-negative integer $k$ such that $\sqrt{2}^k u$ has entries in $\Z[\omega]$.
\end{definition}

By theorem \ref{unitary-iff-clifford+t}, all $U$ that can be implemented by Clifford+T circuits have entries in $\Z[\sq, i]$, so they have a well-defined least denominator exponent.

The least denominator exponent is a measure of how far a unitary matrix is from being implementable by Clifford gates alone, as Clifford gates have entries in $\Z[i]$ and thus have least denominator exponent $0$.
We can use the following theorem to reduce the least denominator exponent of a unitary matrix by applying $H$ and $T$ gates.

\begin{theorem}\label{lde-reduction}
	Let $U \in \mathcal{CT}(2)$,
	let $\text{lde-tl}(U)$ denote $\lde(|z|^2),$ where $z$ is its top left entry.
	Assuming $\text{lde-tl}(U) \geq 4$
	Then there exists $k \in \{0, -1, -2, -3\}$ such that $\text{lde-tl}(H T^k U) = \text{ldt-tl}(U) - 1$.
\end{theorem}

There are only finitely many unit vectors $v$ such that $\lde(|v|^2) \leq 4$.
As the column of a unitary matrix $U \in \mathcal{CT}(2)$ is a unit vector, we can reduce $\text{ldt-tl}(U)$ to a small enough number, after which we can find a Clifford+T circuit that implements it by brute-force search.

We list the exact steps of the algorithm \ref{exact-synthesis-algorithm} for clarity.

\begin{algorithm}[ht]\label{exact-synthesis-algorithm}
	\caption{Exact Synthesis of Single Qubit Unitaries over Clifford+T}
	\hspace*{\algorithmicindent} \textbf{Input:} $U \in \mathcal{CT}(2)$ \\
\hspace*{\algorithmicindent} \textbf{Output:} a sequence of $H$ and $T$ gates that implements $U$
\begin{algorithmic}[1]
	\State $M \gets U, R \gets I$ \Comment{$M$ is temporary matrix, $R$ stores the result}
	\While{$\text{lde-tl}(M) \geq 4$} \label{while-loop-exact-synthesis}
	\State find $k \in \{0, -1, -2, -3\}$ such that $\text{lde-tl}(H T^k M) = \text{lde-tl}(M) - 1$ \Comment{Guaranteed by theorem \ref{lde-reduction}}
\State $M \gets H T^k M$ 
\State $R \gets R T^{-k} H$
\EndWhile 
\State Using brute force to find a sequence of $H$ and $T$ gates that implements $M$. Label this sequence as $M_s$\Comment{There are only finitely many options with $\lde \leq 2$}
\State $R \gets M_s \cdot R$
\State \textbf{return} $R$
\end{algorithmic}
\end{algorithm}

The major work of the algorithm is the while loop in line \ref{while-loop-exact-synthesis}.
After each iteration, the result is appended with $TH, SH, $ or $TSH$. 
So the run time of the algorithm is linear in the number of gates in the output circuit.
Combining with Mastumoto-Amano normal form discussed in the next section, it can be shown that the T-count of this algorithm is optimal.

The case of multi-qubit unitaries that can be generated by Clifford+T circuits is more complicated. 
Kliuchnikov et al hypothesize that any unitary matrix that can be implemented by a Clifford+T circuit has entries in $\Z\left[\frac{1}{\sqrt{2}}, i\right]$.
